For a boundary condition $B$ of a domain $D$ in a Riemannian manifold $M$,
denote by $C^\infty_{c,B}(\bar{D})$ the space of smooth functions on $\bar D$ with compact support
which satisfy $B$ along the boundary~$\partial D$ of $D$.

\begin{Definition}%\label{bcg}
We say that a boundary condition $B$ for a smooth domain $D$ is \emph{non-positive}
if~$u\in C^\infty_{c,B}(\bar D)$ implies $2u\nabla u=\nabla_\nu u^2\le0$ along $\partial D$.
\end{Definition}

\begin{exas}
Dirichlet and Neumann boundary conditions are non-positive.
Robin boundary conditions $\alpha u+\beta\nabla_\nu u=0$ are non-positive if $\alpha\beta>0$.
These kinds of boundary conditions are also self-adjoint and elliptic in the usual sense.
\end{exas}

%%%%%%%%%%%%%%%%%%%%%%%%%%%%%%%%%%%%%%%%%%%%%%
\section{Vanishing of point spectrum}\label{secrell}
%%%%%%%%%%%%%%%%%%%%%%%%%%%%%%%%%%%%%%%%%%%%%%

For a $C^1$ vector field $X$ and a $C^2$ function $u$ on a Riemannian manifold $M$, we have the following identity
\begin{align}
 2\langle\nabla_{\nabla u}X,\nabla u\rangle
 = 2\langle X,\nabla u\rangle\Delta u + |\nabla u|^2\diver X + \diver\big\{2X(u)\nabla u-|\nabla u|^2X\big\}. \label{rell}
\end{align}
In the context of spectral theory, this identity was used by Donnelly--Garofalo,
see~\cite[Lemma~2.1]{DoGa92}, \cite[Proposition 3.1]{DoGa97}, and Xavier, see~\cite[identity~(1)]{Xavier1988},
but it also occurs in the earlier work of Kazdan--Warner~\cite[identity~(8.1)]{KazdanWarner74c}.
Donnelly and Garofalo attribute \eqref{rell} to Rellich~\cite{Rellich43} in the case of Euclidean spaces;
cf.~(3) respectively II on page 62 of~\cite{Rellich43}.
For the convenience of the reader, we give a~short proof of the identity.

\begin{proof}[Proof of \eqref{rell}]
Since $X(u)=\langle\nabla u, X\rangle$ and $|\nabla u|^2=\langle\nabla u,\nabla u\rangle=\nabla u(u)$, we have
\begin{align*}
 2\diver(X(u)\nabla u)
 &= 2\nabla u(X(u)) + 2X(u)\diver\nabla u \\
 &= 2X(\nabla u(u)) + 2[\nabla u,X](u) - 2X(u)\Delta u \\
 &= 2X\bigl(|\nabla u|^2\bigr) + 2(\nabla_{\nabla u}X)(u) - 2(\nabla_X\nabla u)(u) - 2X(u)\Delta u \\
 &= 2X\bigl(|\nabla u|^2\bigr) + 2\langle\nabla_{\nabla u}X,\nabla u\rangle - 2\langle\nabla_X\nabla u,\nabla u\rangle - 2X(u)\Delta u \\
 &= X\bigl(|\nabla u|^2\bigr) + 2\langle\nabla_{\nabla u}X,\nabla u\rangle - 2X(u)\Delta u \\
 &= \diver\bigl(|\nabla u|^2X\bigr) - |\nabla u|^2\diver X + 2\langle\nabla_{\nabla u}X,\nabla u\rangle - 2X(u)\Delta u,
\end{align*}
which is the assertion.
\end{proof}

\begin{Corollary}%\label{codoga}
If $\Delta u=\varphi u$ for a $C^1$ function $\varphi$ on $M$, then
\begin{gather}
 2\langle\nabla_{\nabla u}X,\nabla u\rangle
= \varphi X(u^2) + |\nabla u|^2\diver X + \diver\big\{2X(u)\nabla u-|\nabla u|^2X\big\} \label{doga2} \\
\hphantom{2\langle\nabla_{\nabla u}X,\nabla u\rangle}{}= \bigl(|\nabla u|^2 - \varphi u^2\bigr)\diver X - X(\varphi)u^2 \notag\\
 \hphantom{2\langle\nabla_{\nabla u}X,\nabla u\rangle=}{}+ \diver\big\{2X(u)\nabla u+\bigl(\varphi u^2-|\nabla u|^2\bigr)X\big\}. \label{doga3}
\end{gather}
\end{Corollary}

\begin{proof}
If $\Delta u=\varphi u$, then
\begin{align*}
 2X(u)\Delta u &= 2\varphi X(u)u = \varphi X\bigl(u^2\bigr) = X\bigl(\varphi u^2\bigr) - X(\varphi)u^2 \\
 &= \diver\bigl(\varphi u^2X\bigr) - \varphi u^2\diver X - X(\varphi)u^2.
\end{align*}
Together with \eqref{rell}, this gives \eqref{doga2} and \eqref{doga3}.
\end{proof}
